\documentclass[11pt,a4paper]{article}
\usepackage[utf8]{inputenc}
\usepackage[T1]{fontenc}
\usepackage{amsmath,amssymb}
\usepackage{booktabs}
\usepackage{geometry}
\usepackage{enumitem}
\usepackage{listings}
\usepackage{xcolor}
\usepackage{hyperref}

\geometry{margin=2.5cm}

\lstset{
  basicstyle=\ttfamily\small,
  breaklines=true,
  frame=single,
  columns=fullflexible
}

\title{Empirical Protocol Audit and Contract Binding Analysis\\
for Synchronous C++ gRPC Streams}
\author{Mihai Tira}
\date{November 2026}

\begin{document}

\maketitle

\section{Overview and Evidence Classification}

Unlike stateless unary RPCs, streaming RPCs in gRPC establish stateful, multi-message channels over HTTP/2. In C++, synchronous streaming endpoints (\texttt{grpc::ClientReader}, \texttt{grpc::ClientWriter}, \texttt{grpc::ClientReaderWriter}, and their server counterparts) require callers to obey two classes of protocol invariants:
\begin{enumerate}
    \item \textbf{Level 1 (Call Sequencing and Stream Lifecycle):} Ordering constraints over \texttt{Read()}, \texttt{Write()}, \texttt{WritesDone()}, and \texttt{Finish()}, ensuring that client streams half-close write directions before finalization and that every active stream invokes \texttt{Finish()} before going out of scope.
    \item \textbf{Level 2 (Intra-Message Payload Progression):} State-dependent constraints on Protobuf message fields, such as mandatory \texttt{oneof} initialization variants (e.g., sending configuration metadata on the first message and raw payload chunks on subsequent messages).
\end{enumerate}

To evaluate these invariants systematically across open-source systems, we group empirical evidence into four categories:
\begin{itemize}
    \item \textbf{Category A (Confirmed Historical Defects):} Bugs and pull requests confirmed and fixed by upstream maintainers in repository version history.
    \item \textbf{Category B (Documented Schema Invariants):} Protocol ordering and payload rules specified in production \texttt{.proto} files.
    \item \textbf{Category C (Controlled Fault Injection):} Synthetic protocol mutations used to test analyzer recall and precision.
    \item \textbf{Category D (Candidate Static-Analysis Findings):} Unconfirmed protocol and lifecycle bugs discovered in current repository revisions during our static audit.
\end{itemize}

\section{Audited Open-Source C++ gRPC Repositories}

\subsection{Google Cloud C++ SDK (\texttt{googleapis/google-cloud-cpp})}

\paragraph{Cloud Storage (\texttt{Storage.BidiWriteObject}).}
Defined in \texttt{google/storage/v2/storage.proto}, \texttt{BidiWriteObject} requires the initial \texttt{BidiWriteObjectRequest} on a stream to populate the \texttt{first\_message} \texttt{oneof} (\texttt{write\_object\_spec} or \texttt{upload\_id}). Subsequent writes in the upload loop must clear this specification and send \texttt{checksummed\_data} chunks until the final message sets \texttt{finish\_write = true}, followed by \texttt{WritesDone()} and \texttt{Finish()}. Upstream pull requests (\#16073, \#16083, \#16112, \#16270) document recurring state-tracking fixes around stream initialization and teardown. In synchronous C++ (\texttt{google/cloud/storage/internal/grpc/stub.cc}), the implementation populates \texttt{mutable\_write\_object\_spec()} prior to the chunk loop and calls \texttt{clear\_write\_object\_spec()} after the first \texttt{Write()}.

\paragraph{Cloud Speech-to-Text (\texttt{Speech.StreamingRecognize}).}
In \texttt{cloud\_speech.proto}, \texttt{StreamingRecognizeRequest} uses a \texttt{streaming\_request} \texttt{oneof} where the first message must set \texttt{streaming\_config} (case 1) and all subsequent messages must set \texttt{audio\_content} (case 2):
\[
S_{\text{client}} = \,!\text{Req}[\text{streaming\_config}];\; \mu X.\, \oplus\{\text{audio}: \,!\text{Req}[\text{audio\_content}];\, X,\; \text{done}: \text{WritesDone}();\, ?\text{Resp}^*;\, \text{Finish}()\}
\]
Issue \#9884 documents client failures when audio chunks are sent before configuration or after \texttt{WritesDone()}.

\subsection{WeNet Speech Recognition (\texttt{wenet-e2e/wenet})}

WeNet's \texttt{ASR.Recognize} bidirectional stream (\texttt{runtime/core/grpc/wenet.proto}) uses a \texttt{RequestPayload} \texttt{oneof} requiring \texttt{decode\_config} on the first write and \texttt{audio\_data} on subsequent writes.

In the WebSocket server (\texttt{websocket\_server.cc}), Issue \#813 and PR \#911 fixed a crash where an immediate client disconnect caused \texttt{OnSpeechEnd()} to fail on \texttt{CHECK(feature\_pipeline\_ != nullptr)}. However, the synchronous gRPC server handler (\texttt{GrpcConnectionHandler::operator()} in \texttt{runtime/core/grpc/grpc_server.cc}) remains unguarded:
\begin{lstlisting}[language=C++]
while (stream_->Read(request_.get())) {
  if (!got_start_tag_) {
    OnSpeechStart(); // Initializes feature_pipeline_
  } else {
    OnSpeechData();
  }
}
OnSpeechEnd(); // Calls CHECK(feature_pipeline_ != nullptr) -> SIGABRT on empty stream
\end{lstlisting}
If a client opens a stream and closes it with zero messages, \texttt{Read()} returns \texttt{false} immediately and \texttt{OnSpeechEnd()} aborts the server process (\texttt{SIGABRT}). Furthermore, omitting a \texttt{request\_->has\_decode\_config()} guard on the first message lets an out-of-order \texttt{audio\_data} frame silently drop its audio payload and overwrite \texttt{nbest\_} to \texttt{0} from \texttt{DecodeConfig::default\_instance()}, bypassing the hypothesis limit in \texttt{SerializeResult()}.

\subsection{Apache Arrow Flight (\texttt{apache/arrow})}

Arrow Flight (\texttt{cpp/src/arrow/flight/transport/grpc/grpc\_client.cc}) defines RAII stream wrappers (\texttt{GrpcResultStream} and \texttt{FinishableDataStream}) that call \texttt{TryCancel()} and \texttt{stream\_->Finish()} in their destructors when a caller exits early.

However, \texttt{GrpcClientImpl::ListFlights} (lines 892--914) and \texttt{ListActions} (lines 923--941) use raw \texttt{std::unique\_ptr<grpc::ClientReader<...>>}, while \texttt{AuthenticateInternal} (lines 1067--1091) uses a raw \texttt{ClientReaderWriter}:
\begin{lstlisting}[language=C++]
while (!options.stop_token.IsStopRequested() && stream->Read(&pb_info)) {
  FlightInfo::Data info_data;
  RETURN_NOT_OK(internal::FromProto(pb_info, &info_data));
  flights.emplace_back(std::move(info_data));
}
if (options.stop_token.IsStopRequested()) rpc.context.TryCancel();
RETURN_NOT_OK(options.stop_token.Poll());
return FromGrpcStatus(stream->Finish(), &rpc.context);
\end{lstlisting}
When \texttt{FromProto(pb::FlightInfo)} fails (e.g., \texttt{FlightDescriptor::UNKNOWN} or an invalid \texttt{Location} URI at line 907) or when \texttt{stop\_token} is cancelled (\texttt{Poll()} at line 911 in \texttt{ListFlights} or line 939 in \texttt{ListActions}), \texttt{RETURN\_NOT\_OK} returns early before \texttt{stream->Finish()}. Similarly, while \texttt{AuthenticateInternal} invokes \texttt{Finish()} on its normal exit (line 1089), it omits \texttt{Finish()} when \texttt{auth\_handler\_->Authenticate()} fails (line 1074) and on the \texttt{!finished\_writes} branch (lines 1078--1081). By contrast, \texttt{ListActions}'s \texttt{FromProto(pb::ActionType)} check (line 935) is statically flagged but dynamically unreachable because that overload unconditionally returns \texttt{Status::OK()}. Although \texttt{\~{}ClientContext()} still cancels and reclaims the C-core HTTP/2 call, omitting \texttt{Finish()} prevents \texttt{POST\_RECV\_STATUS} from firing, skipping \texttt{ClientMiddleware::CallCompleted} callbacks and discarding trailing status.

\subsection{BuildGrid BuildBox (\texttt{BuildGrid/buildbox-common})}

In \texttt{common/buildboxcommon\_grpccasclient.cpp}, \texttt{GrpcCASClient::stage} initiates a two-phase \texttt{StageTree} bidirectional stream. When the staged directory is destroyed, \texttt{\~{}GrpcStagedDirectory()} (lines 1317--1323) sends an empty \texttt{StageTreeRequest} and calls \texttt{WritesDone()}, but exits without reading the server's cleanup acknowledgment response or calling \texttt{Finish()}. Although \texttt{buildboxcasd} still finishes directory deletion asynchronously without polling \texttt{ctx->IsCancelled()}, this omission violates \texttt{local\_cas.proto}, aborts the server acknowledgment write via HTTP/2 \texttt{RST\_STREAM}, drops server cleanup errors, bypasses \texttt{POST\_RECV\_STATUS}, and causes immediate re-stage collisions on the same path. Additionally, \texttt{uploadFdRegion} and \texttt{download} throw C++ exceptions on local I/O errors inside active \texttt{Write}/\texttt{Read} loops without calling \texttt{Finish()}.

\subsection{NVIDIA Triton, gRPC Interceptors, and NVIDIA Riva}

\begin{itemize}
    \item \textbf{Triton Inference Server (\texttt{triton-inference-server/client}):} In \texttt{src/c++/library/grpc\_client.cc}, \texttt{StopStream()} calls \texttt{WritesDone()} (or \texttt{TryCancel()}) and joins the reader thread (which calls \texttt{Finish()}), but neither resets \texttt{grpc\_stream\_} nor clears \texttt{enable\_stream\_stats\_}. Subsequent calls to \texttt{AsyncStreamInfer()} invoke \texttt{Write()} on a finished stream: \texttt{libgrpc++} returns \texttt{false} (\texttt{"Stream has been closed."}) rather than crashing, but pushes an unpopped phantom \texttt{RequestTimers} entry that corrupts latency telemetry across later \texttt{StartStream()} sessions. (Calling \texttt{AsyncStreamInfer()} before \texttt{StartStream()} reads uninitialized \texttt{enable\_stream\_stats\_} and dereferences null \texttt{grpc\_stream\_}.)
    \item \textbf{gRPC CachingInterceptor (\texttt{examples/cpp/interceptors/caching\_interceptor.h}):} Intercepts \texttt{keyvaluestore::KeyValueStore::GetValues}; on a cache miss in \texttt{PRE\_SEND\_MESSAGE}, it performs \texttt{stream\_->Write(req)} and \texttt{stream\_->Read(\&resp)} without checking return values, never calls \texttt{stream\_->Finish()}, caches empty strings on failed reads, and unconditionally overwrites \texttt{*status = grpc::Status::OK} in \texttt{PRE\_RECV\_STATUS} (triggering \texttt{RST\_STREAM} teardown and bypassing downstream \texttt{POST\_RECV\_STATUS}). Because \texttt{ClientBidiReactor} callers post \texttt{PRE\_RECV\_STATUS} eagerly at \texttt{StartCall()}, a sound fix must defer \texttt{Finish()} until \texttt{PRE\_SEND\_CLOSE} or stream error.
    \item \textbf{NVIDIA Riva (\texttt{nvidia-riva/cpp-clients}, PR \#56/\#57):} Historical race condition where the writer thread called \texttt{streamer->Finish()} while the reader thread was still executing \texttt{streamer->Read(&response)}.
\end{itemize}

\begin{table}[htbp]
\centering
\small
\begin{tabular}{@{}llll@{}}
\toprule
\textbf{Repository} & \textbf{RPC Method} & \textbf{Level 1 Invariant} & \textbf{Level 2 Invariant} \\
\midrule
\texttt{google-cloud-cpp} & \texttt{Storage.BidiWriteObject} & \texttt{WritesDone} $\to$ \texttt{Finish} & \texttt{write\_object\_spec} on 1st write \\
\texttt{wenet-e2e/wenet} & \texttt{ASR.Recognize} & Non-empty stream before \texttt{OnSpeechEnd} & \texttt{decode\_config} then \texttt{audio\_data} \\
\texttt{apache/arrow} & \texttt{Flight.ListFlights} / \texttt{ListActions} / \texttt{Handshake} & \texttt{Finish()} on early return / cancel & N/A \\
\texttt{buildbox-common} & \texttt{LocalCAS.StageTree} & 2-phase \texttt{Read} ack + \texttt{Finish()} & Stage request then empty cleanup \\
\texttt{triton-client} & \texttt{ModelStreamInfer} & No \texttt{Write()} after \texttt{StopStream()} & N/A \\
\texttt{grpc/grpc} & \texttt{KeyValueStore.GetValues} & Mandatory \texttt{Finish()} before status & N/A \\
\bottomrule
\end{tabular}
\caption{Summary of audited C++ gRPC streaming protocols and invariants.}
\label{tab:audited-repos}
\end{table}

\section{Protocol-to-Compiler Contract Binding}

To connect Protobuf service definitions to Clang AST static analysis without modifying generated \texttt{.grpc.pb.h} files, we compared three contract-binding strategies:
\begin{enumerate}
    \item \textbf{In-Header Attributes:} Injecting \texttt{[[clang::annotate(...)]]} into \texttt{.grpc.pb.h}. While lookup is $O(1)$, regenerating headers with standard \texttt{protoc} silently strips annotations and invalidates build caches.
    \item \textbf{JSON FSM Sidecars (\texttt{.session.json}):} Compiling \texttt{.proto} session options into standalone \texttt{grpc.session.v1} state machine files and matching C++ call sites via Clang AST \texttt{DeclContext} qualified names (\texttt{package::Service::Stub::Method}).
    \item \textbf{Hybrid Binding (Selected):} Using \texttt{.session.json} sidecars as the primary contract representation so unmodified open-source headers can be analyzed directly, with optional attribute annotations reserved for custom wrapper classes.
\end{enumerate}

\subsection{Two-Tier Verification Mapping}

Table~\ref{tab:ast-mapping} summarizes how Clang AST nodes map to transitions in the \texttt{.session.json} state machine.

\begin{table}[htbp]
\centering
\small
\begin{tabular}{@{}lll@{}}
\toprule
\textbf{C++ Stream Operation} & \textbf{Level 1 FSM Transition} & \textbf{Level 2 Payload Check} \\
\midrule
\texttt{stream->Write(msg)} & Advance along \texttt{Write} edge & Check active \texttt{oneof} case on \texttt{msg} \\
\texttt{stream->Read(\&resp)} & Advance along \texttt{Read} edge & Branch on response variant / bool return \\
\texttt{stream->WritesDone()} & Transition to half-closed state & Forbid subsequent \texttt{Write()} calls \\
\texttt{stream->Finish()} & Transition to terminal \texttt{FINISHED} state & Pluck final RPC status \\
Function return / scope exit & Verify state is terminal & Report HTTP/2 stream leak if non-terminal \\
\bottomrule
\end{tabular}
\caption{Mapping between C++ gRPC AST operations and session automaton transitions.}
\label{tab:ast-mapping}
\end{table}

\end{document}
